\documentclass[11pt,a4paper]{article}

\usepackage[utf8]{inputenc}
\usepackage[T1]{fontenc}
\usepackage[ngerman]{babel}
\usepackage{lmodern}
\usepackage{geometry}
\geometry{
    a4paper,
    top=2.2cm,
    bottom=2.2cm,
    left=2.5cm,
    right=2.5cm
}
\usepackage{microtype}
\usepackage{parskip}
\usepackage{titlesec}
\usepackage{booktabs}
\usepackage{tabularx}
\usepackage{enumitem}
\usepackage{amssymb}
\usepackage{hyperref}

\hypersetup{
    colorlinks=true,
    linkcolor=black,
    citecolor=black,
    urlcolor=black
}

\titleformat{\section}{\large\bfseries}{\thesection}{1em}{}
\titleformat{\subsection}{\normalsize\bfseries}{\thesubsection}{1em}{}
\titlespacing*{\section}{0pt}{1.3ex plus 0.3ex minus 0.2ex}{0.5ex}
\titlespacing*{\subsection}{0pt}{1.0ex plus 0.2ex minus 0.2ex}{0.3ex}

\begin{document}

\noindent
{\LARGE\textbf{Abstract und Abschlussbericht: Uni Speedrun}}\\[0.2cm]
{\large Objektorientierte und funktionale Programmierung mit Python}\\[0.3cm]
\textbf{Justus Wächter} \hfill \textbf{Matrikelnummer:} IU14173187\\
\textbf{Portfoliophase 3} \hfill \textbf{Kurs:} DLBDSOOFPP01\_D
\vspace{0.3cm}
\hrule
\vspace{0.4cm}

\section{Ausgangslage und Zielsetzung des Projekts}
Im flexiblen Fernstudienmodell tragen Studierende die Verantwortung für die zeitliche Organisation ihres Studienfortschritts. Während hochschuleigene Portale primär Notenübersichten bereitstellen, fehlt häufig eine dynamische Verlaufs- und Zeitplanung, die Verzögerungen, Bearbeitungszeiten und Korrekturphasen transparent verknüpft. Das Projekt \textit{Uni Speedrun} schließt diese Lücke durch ein leichtgewichtiges, lokal ausführbares Terminal-User-Interface (TUI). 

Ziel der finalen Phase 3 war die vollständige Implementierung des in Phase 1 konzipierten und in Phase 2 architektonisch modellierten Systems. Der Prototyp liefert eine automatisierte Fortschritts- und Zeitprognose, unterstützt einen praxisnahen Modul-Lebenszyklus, persistiert alle Daten in einer lokalen SQLite-Datenbank und setzt die Prinzipien einer sauberen objektorientierten Schichtenarchitektur unter Einbindung des Feedbacks aus Phase 2 konsequent um.

\section{Methodik, Schichtenarchitektur und technische Umsetzung}
Die Umsetzung erfolgte in Python 3.10+ unter strikter Einhaltung der SOLID-Prinzipien. Als Reaktion auf das Feedback der Portfoliophase 2 wurde das Fachmodell erweitert sowie eine dedizierte Controller-Schicht zur klaren Trennung von UI und Persistenz eingeführt:

\subsection{Fachmodell (Domain Layer)}
Das Fachmodell kapselt die betriebswirtschaftliche Logik der Studienorganisation über drei miteinander verknüpfte Entity-Klassen sowie eine Status-Enumeration:
\begin{itemize}[leftmargin=*,noitemsep,topsep=2pt]
    \item \textbf{Prüfungsleistung:} Eigene Entity zur Modellierung von Prüfungen (\texttt{titel}, \texttt{datum}, \texttt{note}, \texttt{bestanden}). Kapselt die Validierung gültiger Notenwerte (1.0 bis 5.0) und entlastet das Modul von reinen Prüfungsverwaltungsaufgaben.
    \item \textbf{Modul:} Repräsentiert ein Studienfach und hält eine optionale Referenz auf eine \texttt{Prüfungsleistung}. Kapselt fachliche Attribute (Name, ECTS, Dauer in Tagen, Reihenfolge, Status) über validierende Python-Properties (\texttt{@property} und Setter), um ein Umgehen fachlicher Regeln auch bei späteren Attributänderungen im Speicher auszuschließen.
    \item \textbf{Studienplan:} Verwaltet die Komposition der Module als geordnete Menge. Er berechnet erreichte ECTS, den prozentualen Fortschritt und über \texttt{prognostiziertes\_studienende()} die Zeitprognose.
    \item \textbf{Modulstatus (Enumeration):} Steuert die Zustandsverwaltung über \texttt{GEPLANT}, \texttt{AKTIV}, \texttt{WARTE\_AUF\_ERGEBNIS} und \texttt{ABGESCHLOSSEN}. Unzulässige Übergänge werden durch Fachmethoden verhindert.
\end{itemize}

\subsection{Anwendungssteuerung (Controller Layer)}
Die Klasse \texttt{DashboardController} koordiniert Benutzeraktionen, Geschäftsabläufe und den Repository-Zugriff. Die Hauptanwendung \texttt{UniSpeedrunApp} fungiert als \textit{Composition Root} (Dependency Injection), erzeugt das Repository und den Controller und übergibt letzteren an die Benutzeroberfläche.

\subsection{Persistenzschicht (Repository Layer)}
Das \textit{Repository Pattern} trennt die Datenhaltung von der Fachlogik. Die Schnittstelle \texttt{StudienplanRepository} (abstrakte Basisklasse) wird von \texttt{SQLiteStudienplanRepository} implementiert (Generalisierung via DIP). Weder Fachmodell noch Benutzeroberfläche enthalten SQL-Befehle.

\subsection{Präsentationsschicht (Textual TUI Layer)}
Die Benutzeroberfläche basiert auf dem Framework \textit{Textual} mit TCSS-Stylesheets (\texttt{app.tcss}) und ist rein auf Darstellung und Eingabeempfang beschränkt: \texttt{DashboardScreen}, \texttt{EinstellungenScreen}, \texttt{ModulBearbeitenScreen} und \texttt{ArchivScreen}.

\newpage

\section{Ergebnisse und Funktionsumfang des Prototyps}
Der entwickelte Prototyp deckt alle Anforderungen der Aufgabenstellung vollständig ab:
\begin{itemize}[leftmargin=*,noitemsep,topsep=2pt]
    \item \textbf{Zentrales Dashboard:} Zeigt Zieldatum, prognostiziertes Studienende und die kalendarische Abweichung (z.\,B. \textit{+4 Monate}). Das aktive Modul wird mit Resttagen hervorgehoben. Eine zentrierte Fortschrittskarte visualisiert den ECTS-Stand und den Fortschrittsbalken (0--100\,\%).
    \item \textbf{Praxisnaher Modul-Lebenszyklus:} Über das Tastenkürzel \texttt{[M]} wird das aktive Modul nach Prüfungsabgabe in den Status \texttt{WARTE\_AUF\_ERGEBNIS} versetzt. Dadurch wird der aktive Bearbeitungsplatz freigegeben und das nächste Modul kann gestartet werden, noch während die Korrektur läuft.
    \item \textbf{Einstellungen und Modulverwaltung:} Über \texttt{[S]} können Studienziel, Ziel-ECTS, Zieldauer und Startdatum im Format \texttt{TT.MM.JJJJ} editiert werden (Auto-Save). Module lassen sich frei anlegen, löschen, bearbeiten und über Pfeiltasten ($\blacktriangle$/$\blacktriangledown$) oder Positionsangabe neu reihen.
    \item \textbf{Archiv und Notenübersicht:} Nach Vorliegen der Bewertung wird das Modul abgeschlossen. Der \texttt{ArchivScreen} (\texttt{[A]}) listet alle Leistungen auf und berechnet automatisch den Notendurchschnitt.
    \item \textbf{Qualitätssicherung:} Die funktionale Korrektheit aller Entitäten, Controller-Abläufe und Persistenzoperationen ist durch \textbf{56 automatisierte Tests} (\texttt{pytest}) überprüft.
\end{itemize}

\section{Kritische Reflexion und Lessons Learned}
Im Rahmen der Implementierung wurden zentrale Entwurfsentscheidungen und Python-Besonderheiten reflektiert:
\begin{enumerate}[leftmargin=*,noitemsep,topsep=2pt]
    \item \textbf{Kapselung \& Property-Validierung in Python:} Da Python keine strikte Zugriffskontrolle (wie \texttt{private}) erzwingt, gewährleisten konsequent eingesetzte Setter-Properties (\texttt{@property}), dass Invarianten (z.\,B. ECTS $> 0$, Note zwischen 1.0 und 5.0) nicht durch nachträgliche Attributzuweisungen verletzt werden können.
    \item \textbf{Entkopplung durch Controller \& DIP:} Durch die Einführung des \texttt{DashboardController}s wurde die Benutzeroberfläche vollständig von Persistenz- und Steuerungslogik befreit. Die Klassen bleiben hochgradig testbar und unabhängig austauschbar.
    \item \textbf{Lebenszyklusabhängigkeit bei Komposition:} Die Bindung von Modulen an den Studienplan wurde fachlich als Komposition realisiert, da Verlaufsdaten (Start-, Prüfungs-, Abschlussdatum) an die individuelle Planung gekoppelt sind.
    \item \textbf{TUI-Rendering und Event-Steuerung:} In Textual wurden Rich-Markup-Konflikte bei Tastatur-Shortcuts (\texttt{[S]}, \texttt{[M]}) durch Deaktivieren des Markups behoben. Modale Dialoge wurden über \texttt{dismiss()} asynchron an die aufrufenden Screens zurückgebunden.
\end{enumerate}

\section{Fazit}
Mit \textit{Uni Speedrun} wurde ein modularer Prototyp für die Studienverlaufsplanung geschaffen. Das Projekt demonstriert objektorientierte Modellierung, automatisierte Tests und eine klare Schichtenarchitektur in Python. Die im Konzept festgelegten Kernfunktionen wurden umgesetzt und anhand der Tests überprüft.

\end{document}
